% Responsable: Arteaga Alava Danela Dayana · Ronda 2
\section{Análisis con el modelo de Kano}

\subsection{Características evaluadas}

Se eligieron 8 características que el socio percibe de forma directa, porque el cuestionario de Kano mide la satisfacción del usuario final. Cada característica se relaciona con un requisito de la Tabla~\ref{tab:requisitos}.

\begin{itemize}
\item \textbf{C1. Ingreso con código QR} (RF-05).
\item \textbf{C2. Aviso de vencimiento} (RF-06).
\item \textbf{C3. Rutina en la aplicación} (RF-08).
\item \textbf{C4. Registro de medidas} (RF-09).
\item \textbf{C5. Reserva de clases grupales} (RF-11).
\item \textbf{C6. Congelación de membresía} (RF-12).
\item \textbf{C7. Comprobante de pago} (RF-03).
\item \textbf{C8. Validación rápida del ingreso} (RC-01).
\end{itemize}

\subsection{Cuestionario}

Cada pregunta se responde con una de estas cinco opciones: (1) me gusta, (2) lo espero, (3) me es indiferente, (4) puedo tolerarlo y (5) me disgusta. Las preguntas impares son funcionales y las pares son disfuncionales. Las preguntas 1 y 2 corresponden a C1, las preguntas 3 y 4 a C2, y así sucesivamente hasta C8.

Las preguntas aplicadas fueron las siguientes:

\begin{enumerate}
\item ¿Cómo se sentiría si pudiera ingresar al gimnasio mostrando un código QR en su celular?
\item ¿Cómo se sentiría si no pudiera ingresar con código QR y tuviera que dictar su número de cédula?
\item ¿Cómo se sentiría si recibiera un correo 5 días antes de que venza su membresía?
\item ¿Cómo se sentiría si no recibiera ningún aviso antes de que venza su membresía?
\item ¿Cómo se sentiría si pudiera ver en su celular la rutina que le toca cada día?
\item ¿Cómo se sentiría si no pudiera ver su rutina en el celular y tuviera que pedirla al entrenador?
\item ¿Cómo se sentiría si el sistema guardara sus medidas corporales y su fecha para ver su avance?
\item ¿Cómo se sentiría si el sistema no guardara sus medidas corporales?
\item ¿Cómo se sentiría si pudiera reservar su cupo en una clase grupal?
\item ¿Cómo se sentiría si no pudiera reservar y tuviera que llegar temprano para alcanzar cupo?
\item ¿Cómo se sentiría si pudiera congelar su membresía entre 7 y 30 días sin perder esos días?
\item ¿Cómo se sentiría si no pudiera congelar su membresía y perdiera los días que no asiste?
\item ¿Cómo se sentiría si cada pago quedara registrado con un número de comprobante?
\item ¿Cómo se sentiría si sus pagos no quedaran registrados con un comprobante?
\item ¿Cómo se sentiría si la validación de su ingreso tardara como máximo 2 segundos?
\item ¿Cómo se sentiría si la validación de su ingreso tardara más de 2 segundos y se formaran filas?
\end{enumerate}

\subsection{Tabla de evaluación}

La Tabla~\ref{tab:evaluacion-kano} cruza la respuesta funcional, con la respuesta disfuncional para obtener la categoría de cada respuesta \cite{kano1984}.

\begin{table}[H]
\centering\footnotesize
\caption{Tabla de evaluación de Kano}\label{tab:evaluacion-kano}
\begin{tabular}{@{}lccccc@{}}
\toprule
\textbf{Funcional $\backslash$ Disfuncional} & \textbf{1. Me gusta} & \textbf{2. Lo espero} & \textbf{3. Indiferente} & \textbf{4. Lo tolero} & \textbf{5. Me disgusta}\\
\midrule
1. Me gusta & Q & A & A & A & O\\
2. Lo espero & R & I & I & I & M\\
3. Indiferente & R & I & I & I & M\\
4. Lo tolero & R & I & I & I & M\\
5. Me disgusta & R & R & R & R & Q\\
\bottomrule
\end{tabular}
\end{table}

A: atractiva. O: unidimensional. M: obligatoria. I: indiferente. R: inversa. Q: cuestionable.
